\newpage
\begin{center}
\textbf{ВВЕДЕНИЕ}
\end{center}
\addcontentsline{toc}{section}{ВВЕДЕНИЕ}

\textbf{Актуальность темы.} В условиях активной цифровизации сферы услуг
автоматизация ключевых бизнес-процессов ресторанного бизнеса становится
важным условием повышения качества обслуживания. Процесс бронирования
столиков остаётся одним из наиболее трудоёмких в части ручного учёта:
администраторы вынуждены вести резервации в журналах или электронных
таблицах, что может приводить к ошибкам, двойным бронированиям и потере
клиентов. Разработка специализированного веб-приложения, автоматизирующего
цикл резервации столиков, является актуальной практической задачей.

\textbf{Степень разработанности проблемы.} На рынке существуют платформы
бронирования и автоматизации ресторанного бизнеса, такие как Restoplace,
Yclients и iiko~\cite{restoplace,yclients,iiko}. Однако подобные решения
могут быть избыточными для небольшого заведения, которому требуется
простая и адаптируемая система управления бронированием. Поэтому
разработка гибкого веб-приложения, легко адаптируемого под нужды
конкретного ресторана, остаётся актуальной задачей.

\textbf{Объект исследования}~--- процесс управления бронированием столиков
в ресторане.

\textbf{Предмет исследования}~--- программные средства и технологии
разработки веб-приложений для автоматизации процесса резервации столиков,
уведомления гостей и формирования отчётности.

\textbf{Цель работы}~--- разработка веб-приложения для автоматизации
процесса бронирования столиков в ресторане, обеспечивающего управление
резервациями, уведомление гостей и формирование отчётности.

Для достижения поставленной цели необходимо решить следующие
\textbf{задачи}:
\begin{enumerate}
    \item Провести обзор существующих систем бронирования и анализ
          инструментов разработки с целью обоснования выбора
          технологического стека;
    \item Проанализировать целевую аудиторию и описать функциональные
          требования к приложению посредством диаграммы вариантов
          использования и пользовательских историй (User Story);
    \item Спроектировать модель данных: разработать диаграмму
          сущность-связь (ER-диаграмму), логическую и физическую схемы
          базы данных;
    \item Разработать макеты страниц (Wireframe) и реализовать базовую
          структуру приложения с вёрсткой шаблонов страниц;
    \item Реализовать систему аутентификации пользователей и интерфейс
          управления данными (CRUD) для столиков и бронирований,
          включая загрузку изображений;
    \item Реализовать систему проверки доступности столиков и управления
          временем бронирования с поддержкой экспорта отчётов о
          занятости (CSV, PDF);
    \item Разработать систему уведомлений и генерации
          PDF-подтверждений бронирования с отправкой гостям по
          электронной почте;
    \item Создать репозиторий с исходным кодом проекта и выполнить
          развёртывание приложения на хостинге.
\end{enumerate}

\textbf{Методы исследования:} анализ и сравнение программных продуктов,
проектирование архитектуры веб-приложений, объектно-ориентированное
программирование, прототипирование интерфейсов.

\textbf{Практическая значимость} работы состоит в том, что разработанное
приложение может быть непосредственно внедрено в работу ресторана и
заменить ручной учёт резервации столиков, сократив количество ошибок и
повысив качество обслуживания гостей.
